% from ega4-17-fr.tex
\oldpage[IV-4]{65}
$k$-algèbre \emph{finie}, composée directe de $k$-algèbres finies locales. Pour exprimer que $\Delta_f$ est une immersion ouverte, on peut donc se borner au cas où $A$ est une $k$-algèbre locale finie, $X=\Spec(A)$ étant donc réduit à un seul point; le corps résiduel de $A$, étant une extension finie de $k$, est nécessairement identique à $k$, et par suite (I, 3.4.9) $X\times_k X$ est réduit à un seul point et $\Delta_f$ est donc nécessairement un \emph{isomorphisme}. Or, puisque $A$ est une $k$-algèbre, l'homomorphisme canonique $A\otimes_k A\to A$ ne peut être bijectif que si $A=k$. C.Q.F.D.

\begin{remark}[17.4.1.2]
\label{IV.17.4.1.2}
Supposons seulement que $f$ soit \emph{localement de type fini}. Alors $\Omega^1_{X/Y}$ est encore un $\sh{O}_X$-Module de type fini (16.3.9), et $\Delta_f$ un morphisme localement de présentation finie (I, 4.3.1). Toute la démonstration de (17.4.1) est donc valable, à condition de remplacer a) par: \emph{la restriction de $f$ à un voisinage convenable de $x$ est un morphisme formellement non ramifié}. On voit en outre que dans ce cas la restriction de $f$ à un voisinage convenable de $x$ est un morphisme \emph{localement quasi-fini}.
\end{remark}

\begin{corollary}[17.4.2]
\label{IV.17.4.2}
Soit $f:X\to Y$ un morphisme localement de présentation finie. Les propriétés suivantes sont équivalentes:
\begin{enumerate}
  \item[a)] $f$ est non ramifié.
  \item[b)] Le morphisme diagonal $\Delta_f:X\to X\times_Y X$ est une immersion ouverte.
  \item[b')] Pour tout morphisme $Y'\to Y$, toute $Y'$-section de $X'=X\times_Y Y'$ est une immersion ouverte.
  \item[c)] On a $\Omega^1_{X/Y}=0$.
  \item[d)] Pour tout $y\in Y$, le $k(y)$-préschéma $f^{-1}(y)$ est non ramifié sur $k(y)$.
  \item[d')] Pour tout $y\in Y$, le $k(y)$-préschéma $f^{-1}(y)$ est isomorphe à un préschéma de la forme $\coprod_{\lambda\in L}\Spec(K_\lambda)$, où, pour tout $\lambda\in L$, $K_\lambda$ est une extension finie et séparable de $k(y)$.
  \item[e)] Pour tout $x\in X$, l'anneau $\sh{O}_{X,x}$ est une $\sh{O}_{Y,f(x)}$-algèbre formellement non ramifiée pour les topologies discrètes.
\end{enumerate}
\end{corollary}

\begin{corollary}[17.4.3]
\label{IV.17.4.3}
Si $f:X\to Y$ est non ramifié, alors $f$ est localement quasi-fini (Err$_{\mathrm{III}}$, 20).
\end{corollary}

\begin{proposition}[17.4.4]
\label{IV.17.4.4}
Soient $Y$ un préschéma localement noethérien, $f:X\to Y$ un morphisme localement de type fini, $x$ un point de $X$, $y=f(x)$. Posons $A=\sh{O}_{Y,y}$, $B=\sh{O}_{X,x}$, qui sont des anneaux locaux noethériens, et soit $k$ le corps résiduel de $A$. Alors les conditions équivalentes a) à e) du théorème (17.4.1) sont aussi équivalentes à chacune des suivantes:
\begin{enumerate}
  \item[f)] $\widehat{B}\otimes_{\widehat{A}}k$ est un corps, extension finie séparable de $k$ (ce qui entraîne que $\widehat{B}$ est une $\widehat{A}$-algèbre finie).
  \item[f')] $\widehat{B}$ est une $\widehat{A}$-algèbre formellement non ramifiée pour les topologies adiques.
\end{enumerate}
Si de plus $k(x)=k(y)$, ou si $k$ est séparablement clos, ces conditions équivalent aussi à:
\begin{enumerate}
  \item[f'')] L'homomorphisme $\widehat{A}\to\widehat{B}$ est surjectif.
\end{enumerate}
\end{proposition}

Notons d'abord, par le même raisonnement que dans (0, 19.3.6), qu'il revient au même de dire que $B$ est une $A$-algèbre formellement non ramifiée pour les topologies préadiques, ou que $\widehat{B}$ est une $\widehat{A}$-algèbre formellement non ramifiée pour les topologies adiques. D'autre part, l'hypothèse que $f$ est localement de type fini entraîne que $\Omega^1_{B/A}$

